\subsection{Glyph: \glyph{Catalysis}}
\label{sec:catalysis}

A catalysis is a particular case of stimulation, where the effector affects positively the flux of a process represented by the target process. The positive effect on the process is due to the lowering of the activation energy of a reaction.

\begin{glyphDescription}

\glyphSboTerm
SBO:0000172 ! catalysis

\glyphOrigin
One \glyph{EPN} (\sect{EPNs}) or \glyph{logical operator} (\sect{logic}).

\glyphTarget
One \glyph{process node} (\sect{PNs}).

\glyphSymbol
The target extremity of a \glyph{catalysis} carries an empty circle, as shown in \fig{catalysis}.

\end{glyphDescription}

\begin{figure}[H]
  \centering
  \includegraphics{images/build/catalysis.pdf}
  \caption{The \PD glyph for \glyph{catalysis}.}
  \label{fig:catalysis}
\end{figure}
